Fixed-rate bond valuation requires consistency between payment timing, discounting and the interest-rate variable being changed. This dissertation examines calculations for regular bonds, smooth term-structure representation and deterministic curve risk. Nominal coupon-compounded valuation is distinguished from discounting with continuously compounded zero rates. Analytical differentiation establishes duration and convexity; a six-parameter Nelson--Siegel--Svensson function represents illustrative maturity-rate observations. Independent cash-flow calculations, high-precision exponential references, finite differences and a partition-of-unity proof assess the implementation. All 234 deterministic tests pass. The illustrative 2030 bond has a clean price of 97.876838 per 100 face value and modified duration of 4.231518 years. The curve-fit root mean squared error is 0.782762 basis points. For the longer risk bond, parallel sensitivity is 0.1697264000 currency units per one-basis-point bump; summed key-rate sensitivity differs by approximately minus \(2.93\times10^{-8}\) currency units, consistent with a cubic correction. The principal limitation is that fitted illustrative rates are assigned a zero-rate interpretation without inference from instrument prices. The evidence supports mathematical consistency under explicit conventions; empirical market validity remains untested.
